\documentclass[sigconf,anonymous,review]{acmart}
\setkeys{acmart.cls}{balance=false}
\usepackage{mathtools}
\usepackage{booktabs}
\usepackage{graphicx}
\setcopyright{none}
\settopmatter{printacmref=false,printfolios=true}
\acmConference[WWW]{The Web Conference}{}{}
\acmDOI{}
\acmISBN{}
\renewcommand\footnotetextcopyrightpermission[1]{}
\setlength{\emergencystretch}{1.5em}
\makeatletter
\AtBeginDocument{%
  \fancyhead[LE]{\ACM@linecountL\@headfootfont The Web Conference}%
  \fancyhead[RO]{\@headfootfont The Web Conference\ACM@linecountR}}
\makeatother
\numberwithin{equation}{section}
\newtheorem{theorem}{Theorem}[section]
\newtheorem{proposition}[theorem]{Proposition}
\newtheorem{lemma}[theorem]{Lemma}
\newcommand{\E}{\mathbb{E}}
\newcommand{\Qmaj}{Q_{\mathrm{maj}}}
\newcommand{\Gs}{G^{*}}
\newcommand{\KY}{K^{\mathrm{Y}}}
\newcommand{\QY}{Q^{\mathrm{Y}}}
\newcommand{\Rset}{\mathcal{R}}
\newcommand{\Bin}{\mathrm{Bin}}
\newcommand{\one}{\mathbf{1}}

\title[Why Hedging Breaks Yuma]{Why Hedging Breaks Yuma: Report-Only Rewards for Decentralized AI Evaluation}
\author{Anonymous Author(s)}

\begin{document}

\begin{abstract}
Bittensor is an open network that pays people, in its own token, to run AI
models (miners) and to test and score them (validators). The network cannot
check which model is actually better, so it pays each validator for agreeing
with the other validators and pays each miner according to the scores.
Nothing stops a validator from also owning a miner, and a validator whose
tests favor a rival may then score its own miner higher. In a model with two
competing miners, we ask how much of a miner a validator can own before
honest scoring stops paying, and how much of the miners' pay a rule that sees
only the validators' reports can then send to the better miner. With a fixed
budget, the best rule follows the majority of reports, and once ownership
passes a threshold it lets the majority decide only part of the miners' pay
and splits the rest evenly. If each validator must put all its weight on one
miner and past scores do not count, Bittensor's rule, Yuma consensus, does as
well as the best rule below the threshold. Real validators split their
weight, and Yuma then rewards hedging: putting a little weight on the miner
one believes is worse. With three validators whose tests are right three times
in four, honest scoring survives ownership of up to 22\% of a miner when
weights are all-or-nothing, but once weights can be split, owning just 1\%
makes a small hedge pay. In simulations that reproduce Bittensor's reward code
exactly, with 11 or 31 validators, a validator that owns nothing earns up to
4.8\% more by blending a tenth of an even split into its weights, and the best
miner's share of pay falls. Yuma's default memory of past scores makes matters
worse: in the three-validator example, honest weights that lean 60/40 or more
toward the miner a validator's test favors are not an equilibrium, even when
no validator owns a miner.
\end{abstract}

\begin{CCSXML}
<ccs2012>
<concept><concept_id>10003752.10010070.10010099</concept_id>
<concept_desc>Theory of computation~Algorithmic game theory and mechanism design</concept_desc>
<concept_significance>500</concept_significance></concept>
<concept><concept_id>10010405.10010455.10010457</concept_id>
<concept_desc>Applied computing~Electronic commerce</concept_desc>
<concept_significance>300</concept_significance></concept>
</ccs2012>
\end{CCSXML}
\ccsdesc[500]{Theory of computation~Algorithmic game theory and mechanism design}
\ccsdesc[300]{Applied computing~Electronic commerce}
\keywords{peer prediction, decentralized AI, Bittensor, Yuma consensus, conflicts of interest, mechanism design}

\maketitle

\begin{quote}
\small\itshape
Sed quis custodiet ipsos custodes?\par
But who will guard the guards themselves?\par
\upshape\hfill -- Juvenal, \emph{Satires} VI
\end{quote}

\section{Introduction}
\label{sec:intro}

Bittensor is an open network that pays, in its own token TAO, for digital
services such as running AI models~\cite{bittensor2026network}. It is divided
into subnets, each an open market for one task, such as answering questions
with a language model. Anyone can join a subnet as a \emph{miner}, who runs a
model that performs the task, or, after locking up enough tokens as
\emph{stake}, as a \emph{validator}, who sends the miners test queries and
scores their answers by splitting one unit of \emph{weight} among them. Every
round the subnet mints a fixed number of new tokens. The miners' part is
divided according to the validators' weights. The validators' part is divided
according to how well each validator's weights agree with those of the others,
because agreement is all the network can check: it never learns which model is
actually better. Figure~\ref{fig:round} sketches one round.

\begin{figure*}[t]
\centering
\includegraphics[width=\textwidth]{figures/round_schematic.pdf}
\caption{One round of a subnet with two miners and three validators. Each
validator splits one unit of weight between the miners; Yuma finds the median
weight on each miner, cuts every weight down to it, and pays the miners by the
cut weights and each validator by the part of its weights that survives.
Carol owns part of Miner~1, and her tests favor Miner~0. Alice and Bob
disagree, so Carol's 20\% hedge moves a fifth of the miners' pay to her own
miner while her pay stays at one half of the validators' pay; in other rounds
a hedge can cost her (Section~\ref{sec:intro-results}). Numbers are shares of
the miners' and of the validators' pay.}
\label{fig:round}
\end{figure*}

Picture a subnet that pays for answers to programming questions. Three
validators, Alice, Bob and Carol, hold equal stakes and score the two leading
miners. Carol also owns a share of one of these two miners, under a different
name. This round her tests say that the rival writes better code, and she must
submit her weights before she sees anyone else's. Suppose she reports what she
found. If Alice and Bob agree with each other, her weights do not change the
miners' pay. If they disagree, her weights break the tie, and her own miner
loses. Suppose she backs her own miner. It then wins whenever the other two
disagree, but her weights clash with the majority more often, and she gives up
part of her validator pay. Carol reports honestly only if the reward for
agreeing outweighs what her miner stands to lose, and the larger her share of
the miner, the more agreement has to be worth.

The same tension arises wherever a platform pays graders without observing the
truth. Crowdsourced labeling and peer grading reward graders for agreeing with
one another~\cite{miller2005peer,prelec2004bayesian,dasgupta2013crowdsourced,
ho2015incentivizing}, public leaderboards of language models are built from
votes on which of two answers is better~\cite{zheng2023judging,
chiang2024chatbot}, and graders may care about the outcome they
grade~\cite{freeman2017crowdsourced,goel2020peer}. Bittensor's reward rule,
Yuma consensus~\cite{bittensor2026yuma}, was built against a different threat.
It cuts every weight down to the stake-weighted median of the weights on the
same miner, a step called \emph{clipping}, so that validators with a minority
of the stake cannot raise a miner's pay on their own. Its documentation calls
the design ``adversarially-resilient when majority stake is honest.'' Whether
scoring honestly is in
each validator's own interest once validators own miners is a different
question, and it is the one we study.

We model a round with two miners and a reward rule that sees only the
validators' weights, which we call \emph{reports}; putting all weight on one
miner amounts to reporting which miner is better. The rule cannot check which
miner is better, it cannot fine anyone, and it pays out fixed budgets to the
miners and to the validators. Each validator can pay a cost to test the
miners, and a test names the better miner with some accuracy. A validator may
own part of a miner, and the rule does not know who owns what. \emph{Honest
scoring} means testing and reporting what the test found; we say it survives,
or is an equilibrium, if no validator gains by deviating while the others
score honestly. We ask how large a share of a miner validators can own before
honest scoring stops surviving, and how much of the miners' pay a rule can then
still send to the better miner. We call the expected fraction of the miners'
pay that reaches the better miner the rule's \emph{quality}. Splitting the pay
evenly gives quality $1/2$, and always paying the better miner would give
quality $1$.

\subsection{What we find}
\label{sec:intro-results}

\paragraph{The best any rule can do.} A validator's temptation to misreport is
its ownership share times how much its report moves the miners' pay. The only
counterweight is its \emph{evaluation premium}, the extra pay it expects for
testing and reporting honestly, and a fixed budget caps that premium. A rule
can therefore tolerate larger holdings only by letting each report move less
money. The best rule does exactly this. It pays the miners according to the
majority verdict and shares the validators' pay among those in the majority.
Up to a threshold ownership share, that is all it needs. Beyond the threshold,
it lets the verdict decide only part of the miners' pay and splits the rest
evenly, and that part shrinks as the largest possible holding grows. No rule
that sees only reports and treats everyone alike does better
(Theorem~\ref{thm:frontier}). In the example, suppose each test is right three
times out of four, testing is free, and the validators' pot equals the miners'
pot. Validators can then own up to 22\% of a miner at no cost in quality
(Figure~\ref{fig:frontier}). Even when a validator may own a whole miner, the
best rule still sends 58\% of the miners' pay to the better miner, against
84\% when nobody owns anything; 84\% is how often the majority of three such
tests is right.

\begin{figure}[t]
\centering
\includegraphics[width=\columnwidth]{figures/frontier.pdf}
\caption{Quality, the share of the miners' pay reaching the better miner,
against the largest share of a miner a validator may own (three validators,
tests right three times in four, free testing, equal pots). Solid: the best any
rule can do (Theorem~\ref{thm:frontier}). Yuma with all-or-nothing weights
matches it up to $2/9\approx22\%$ (circle) and fails beyond. With split weights,
fully decisive honest weights survive only up to $1/108\approx1\%$; the dashed
curve follows the most decisive honest weights that still survive
(Proposition~\ref{prop:soft}), and its quality falls to one half at
$8/51\approx16\%$.}
\label{fig:frontier}
\end{figure}

\paragraph{Where Yuma falls short: hedging.} Yuma also carries past scores into
current pay, a memory we return to below. Without this memory, Yuma does as
well as the best rule below the threshold, provided every validator puts all
its weight on one miner. Real validators split their weight, for instance 80\%
on one miner and 20\% on the other, and Yuma then pays for partial agreement.
An honest validator leans toward the miner its test favors; how far it leans,
from 100/0 down to nearly 50/50, is a convention that validators follow, and we
ask which conventions survive. Split weights also make room for
\emph{hedging}, putting a little weight on the miner one's own test says is
worse. Suppose Carol moves a little weight to her own miner. If Alice and Bob
both side with her test, the hedge costs her a little agreement pay. If both
side against it, the hedged weight earns agreement pay that an all-or-nothing
report would miss. If they split, the hedge moves as large a share of the
miners' pay to her own miner as the weight she shifted, at no cost to her own
pay, as in Figure~\ref{fig:round}. With tests right three times out of four,
the first two effects nearly cancel, and the third tips the balance at a tiny
ownership share. If weights must be all-or-nothing, a validator owning up to
22\% of a miner still does best by scoring honestly; once weights can be split,
owning just 1\% makes a small hedge pay (Theorem~\ref{thm:hedging}). With more
validators the two thresholds move closer together. With tests right only
three times out of five, the second effect wins outright, and hedging pays even
a validator that owns nothing, for every number of validators up to 13
(Table~\ref{tab:thresholds}). Honest weights that lean less, such as 70/30,
survive larger holdings, but the more even the weights, the less of the pay
reaches the better miner, down to one half (Figure~\ref{fig:frontier}).

\paragraph{The deployed code.} To check that these effects survive the details
of the real system, we reimplemented Bittensor's reward
code~\cite{subtensor2026} and checked it against the original bit for bit. It
reproduces the 1\% threshold of the example. We then let one validator hedge in subnets with 11 or 31 validators, up to ten
miners, and equal or concentrated stake, with memory switched off and noisy
test scores; every validator starts from weights that tilt smoothly toward the
miners that test better. In every configuration, a validator that owns nothing
earns more by blending a tenth of an even split into its weights, up to 4.8\%
of its pay, and the best miner's share of pay falls by up to 1.6 percentage
points (Section~\ref{sec:replica}).

\paragraph{Memory makes it worse.} Yuma pays validators through \emph{bonds}, a
running average of each validator's past weights, so weight given today keeps
paying in later rounds. In the example the better miner is drawn afresh each
round, so a bond on yesterday's favorite is a stake in a miner that is no more
likely to be better today: old bonds act like ownership, and honest weights
that lean hard build such bonds quickly. Under Bittensor's default memory, and
with validators who value future rounds nearly as much as the present one,
honest weights that lean 60/40 or more are not an equilibrium in the example,
even when nobody owns a miner. Weights of 51/49 are, when holdings and testing
costs are small, and so is an outcome in which nobody tests at all
(Theorem~\ref{thm:default}). The 51/49 equilibrium sends only 50.6\% of the
miners' pay to the better miner.

\paragraph{Unequal stake and persistent winners.}
Section~\ref{sec:extensions} extends the benchmark in two directions. With
unequal stake, a validator with more stake moves the pay more and is tempted
more, so each validator needs its own premium; we find exactly which premiums
a fixed budget can deliver. Yuma-style pay can even reverse the incentive: a
validator holding 60\% of the stake decides the majority alone, keeps more of
the validators' pay when fewer validators side with it, and so earns more by
ignoring its test, even when it owns nothing. When the better miner tends to
stay better from round to round, past reports already reveal much of what a
fresh test would. In the example, no rule that uses only the current round's
reports can pay for costly testing once the better miner stays better with
probability of about 76\% or more.

\paragraph{A repair.} For the two-miner setting the analysis suggests a simple
repair: pay the miners and the validators according to the miner each
validator weights most, share the validators' pay among the majority, and damp
the miners' pay toward an even split according to the largest holding the
designer allows for.

\subsection{Related work}

\paragraph{Peer prediction.} Peer prediction rewards reports without
observing the truth~\cite{miller2005peer,prelec2004bayesian,witkowski2012peer,
shnayder2016informed,kong2019information}; Faltings and
Radanovic~\cite{faltings2017game} survey the area. Rewards must also cover
evaluation effort~\cite{dasgupta2013crowdsourced,liu2017sequential,
ho2015incentivizing}. Rewarding plain agreement makes graders report what
they expect others to say~\cite{waggoner2014output}, and graders may
coordinate on reports that carry no information, in theory and in
experiments~\cite{jurca2009mechanisms,gao2014trick,gao2019incentivizing};
Theorem~\ref{thm:default}(iii) finds such an outcome in Yuma. Our cap on the
evaluation premium (Lemma~\ref{lem:rent}) asks how a fixed budget can best
reward testing.

\paragraph{Graders who care about the outcome.} Closest to our question, Goel,
Filos-Ratsikas and Faltings~\cite{goel2020peer} find conditions under which
rewarding agreement still yields truthful reports when agents care about the
aggregate of the reports, and Freeman, Lahaie and
Pennock~\cite{freeman2017crowdsourced} fund an agreement reward from market
fees so that arbiters who hold positions in a prediction market vote
truthfully. In both, the reward for reporting is scaled up until it outweighs
the outside incentive. Our budget is fixed, so past an ownership threshold the
split of the miners' pay must give up quality. We find the best quality any
fixed-budget rule can reach, choosing jointly how to pay the validators and
how to split the miners' pay, and a rule that reaches it
(Theorem~\ref{thm:frontier}); Yuma is one rule we measure against this
benchmark.
Scoring rules for self-interested experts and decision
markets~\cite{boutilier2012eliciting,chen2011decision} score reports against
the realized outcome, which a report-only rule never observes.

\paragraph{Committees and robust aggregation.} In committee models with
costly information~\cite{persico2004committee,gerardi2008information,
gershkov2009optimal}, members share the committee's goal, and the chance of
breaking a tie makes information valuable~\cite{austensmith1996information}.
Here agreement pays for testing, while ownership makes a validator care about
the rounds in which its report breaks a tie. Clipping at the median resembles
robust statistical aggregation~\cite{huber1964robust,yin2018byzantine}, and
median voting is strategy-proof under single-peaked preferences when voters
care only about the outcome~\cite{moulin1980strategy}; Yuma also pays
validators for agreeing. These safeguards protect the result from outliers,
and none of them pays anyone to test.

\paragraph{Oracles and Bittensor.} Blockchain oracles share the structure: they
aggregate reports about an external fact that the chain cannot
observe~\cite{eskandari2021sok}. On Bittensor itself, Lui and
Sun~\cite{lui2026bittensor} study the protocol with data and simulation, and the
documentation~\cite{bittensor2026yuma} analyzes groups of large stakeholders
acting together; neither treats validators' holdings in miners or hedging with
split weights. We
study single validators throughout; groups that coordinate raise separate
questions.

\section{Model}
\label{sec:model}

\paragraph{Validators, miners and evaluations.} A round has $n=2m+1\ge3$
validators with equal stake and two miners, labeled $0$ and $1$;
Section~\ref{sec:extensions} allows unequal stake. One miner serves better than
the other. The state $\theta\in\{0,1\}$ names the better miner, and each value
has probability $1/2$. A validator may pay a cost $\eta\ge0$ to evaluate the
miners. An evaluation returns a private signal $s_i\in\{0,1\}$ that names the
better miner with probability $p\in(1/2,1)$, independently across validators
given $\theta$; we write $q=1-p$ and $d=2p-1$. A validator that does not
evaluate learns nothing about $\theta$. Two miners represent a pairwise
quality comparison, as in language-model arenas~\cite{chiang2024chatbot}.
This interpretation retains the binary-signal assumption; it does not
decompose a many-miner subnet into independent pairs. The experiments in
Section~\ref{sec:replica} retain the full normalization across miners.

\paragraph{Reports and rules.} Each validator reports $r_i\in\{0,1\}$, the
miner it names as better. A reward rule observes the reports and nothing else:
not $\theta$, not the signals, not who evaluated. We consider rules that treat
validators alike, so that pay depends on the reports only through the number
$k$ of validators who report $1$ and on a validator's own report. Such a rule is
a pair $(a,v)$. The allocation $a_k\in[0,1]$ is the share of the miner budget
$M>0$ paid to miner $1$ when $k$ validators report $1$; miner $0$ receives
$1-a_k$. The payment $v_{r,k}\ge0$ goes to each validator who reports $r$ when
$k$ validators report $1$. Every round pays out the validator budget $V>0$
exactly:
\begin{equation}\label{eq:budget}
k\,v_{1,k}+(n-k)\,v_{0,k}=V,\qquad k=0,\dots,n.
\end{equation}
In particular, a rule cannot fine a validator or pay out more than $V$. We also
ask that the rule treat the two miners alike, $a_{n-k}=1-a_k$ and
$v_{0,k}=v_{1,n-k}$, and that more support never lower a miner's pay,
$a_k\le a_{k+1}$. We write $\Rset$ for the set of rules with these properties.

\paragraph{Ownership.} A validator may own part of a miner's revenue.
Validator $i$ with ownership $\beta_i$ receives $\beta_iMa_k$ on top of its
validator payment: $\beta_i$ is its share of miner $1$'s revenue minus its share
of miner $0$'s revenue, and the rest of its miner income does not depend on the
reports, because the two miners' revenues sum to $M$. Ownership lies in
$[-\omega,\omega]$ for a bound $\omega\ge0$ known to the designer of the rule;
the rule does not observe $\beta_i$. A validator's payoff in the round is
\begin{equation}\label{eq:payoff}
u_i=v_{r_i,k}+\beta_iMa_k-\eta\,e_i,
\end{equation}
where $e_i=1$ if validator $i$ evaluates and $e_i=0$ otherwise.

\paragraph{Honest evaluation.} Under \emph{honest evaluation}, every validator
evaluates and reports its signal. A rule \emph{supports honest evaluation} at
cost $\eta$ and ownership bound $\omega$ if honest evaluation is an equilibrium
for every ownership profile in $[-\omega,\omega]^n$: when the others evaluate
honestly, no validator gains, whatever its ownership, by skipping the
evaluation, by reporting against its signal, or by randomizing. The definition
asks for existence; other equilibria may coexist with honest evaluation.

\paragraph{Repeated rounds.} Given the exogenous state path, signals are
independent across validators and rounds, with accuracy $p$. Validators
observe only purchased signals and public post-round reports, use independent
private randomization, and discount by $\delta\in[0,1)$.

\paragraph{Quality.} The \emph{quality} of a rule is the expected share of the
miner budget paid to the better miner under honest evaluation,
\begin{equation}\label{eq:quality}
Q(a)=\E\bigl[a_X\bigr],\qquad X\sim\Bin(n,p),
\end{equation}
where $X$ counts the validators whose signal names miner $1$ when $\theta=1$; by
symmetry the same value obtains when $\theta=0$. An even split has quality
$1/2$. The \emph{majority accuracy} $\Qmaj=\Pr[X>n/2]$ is the quality of the
majority rule $a_k=\one\{k>m\}$: the probability that most validators name the
better miner.

\paragraph{Influence and the evaluation premium.} The \emph{influence} of a
rule is the expected change in miner $1$'s share when one validator switches
its report from $0$ to $1$ while the other $n-1$ validators report honest
signals and $\theta=1$:
\begin{equation}\label{eq:influence}
H(a)=\sum_{j=0}^{n-1}\binom{n-1}{j}p^jq^{n-1-j}\,(a_{j+1}-a_j).
\end{equation}
Because the rule treats the miners alike, the same expected change obtains
conditional on the switching validator's own signal. The influence of the
majority rule is the probability that the other $n-1$ signals split evenly, the
pivot probability of committee
models~\cite{austensmith1996information,persico2004committee},
\begin{equation}\label{eq:pivot}
b=\binom{2m}{m}(pq)^m.
\end{equation}
The \emph{evaluation premium} of a rule is what a validator gains in
expectation by evaluating and reporting honestly, over reporting a fixed miner
without evaluating, when the others evaluate honestly:
\begin{equation}\label{eq:premium}
G(v)=\frac Vn-\E\bigl[v_{1,1+J}\bigr],
\end{equation}
where $J$ counts the other validators who report $1$, with $\theta$ equally
likely to be $0$ or $1$. The first term is the expected payment under honest
evaluation, which~\eqref{eq:budget} and the equal treatment of validators fix at
$V/n$ for every rule in $\Rset$. By symmetry, reporting miner $0$ without
evaluating earns the same as reporting miner $1$.

\paragraph{Two rules.} For $\lambda\in[0,1]$, the \emph{damped majority rule}
of strength $\lambda$ pays miner $1$ the share
\begin{equation}\label{eq:damped}
a_k=\frac{1-\lambda}{2}+\lambda\,\one\{k>m\},
\end{equation}
so that a fraction $\lambda$ of the miner budget follows the majority and the
rest is split evenly. \emph{Majority sharing} pays the validator budget in equal
parts to the validators whose report agrees with the majority, and nothing to
the others; when all reports agree, every validator receives $V/n$.

\paragraph{Running example.} Three validators ($m=1$), Alice, Bob and Carol
of the introduction, evaluations that name
the better miner with probability $p=3/4$, and equal budgets $V=M$. Then
$\Qmaj=27/32$ and $b=3/8$. The undamped majority rule ($\lambda=1$) with
majority sharing has influence $H=3/8$ and evaluation premium $G=V/24$.

\section{The frontier}
\label{sec:frontier}

The frontier rests on one incentive constraint and two bounds. The proofs of
the three lemmas are in Appendix~\ref{app:frontier}.

\begin{lemma}[Incentive constraint]\label{lem:ic}
A rule $(a,v)\in\Rset$ supports honest evaluation at cost $\eta$ and ownership
bound $\omega$ if and only if
\[
\eta+\tfrac12\,\omega M H(a)\le G(v).
\]
\end{lemma}

The left side is what honesty costs a validator: the evaluation cost and, for
the worst sign of ownership, the gain from reporting its own miner blindly,
which raises that miner's expected share by $H(a)/2$. The right side is what
honesty earns. Reporting against one's signal after evaluating is never more
tempting than reporting blindly, because it forgoes twice the premium while
moving the allocation by $H(a)$.

\begin{lemma}[Premium cap]\label{lem:rent}
Every rule $(a,v)\in\Rset$ has $G(v)\le\Gs$, where
\[
\Gs=\frac{Vd(\Qmaj-p)}{2npq},
\]
and majority sharing attains $G(v)=\Gs$.
\end{lemma}

Since every honest validator expects $V/n$, a rule raises the premium only by
lowering what a blind report earns. Majority sharing lowers it most: a blind
report lands in the minority more often than an informed one, and the minority
receives nothing.

\begin{lemma}[Quality per unit of influence]\label{lem:influence}
Every allocation $a$ of a rule in $\Rset$ satisfies $Q(a)\le\Qmaj$ and
\[
Q(a)-\frac12\le\frac{\Qmaj-1/2}{b}\,H(a).
\]
\end{lemma}

The damped majority rule~\eqref{eq:damped} meets both bounds: its influence is
$\lambda b$ and its quality is $\frac12+\lambda(\Qmaj-\frac12)$. The proof writes
an allocation as a mixture of symmetric threshold rules and shows that the
majority threshold earns the most quality per unit of influence.

\begin{theorem}[Frontier]\label{thm:frontier}
Let $0\le\eta\le\Gs$. For $\omega>0$, the largest quality of a rule in $\Rset$
that supports honest evaluation at cost $\eta$ and ownership bound $\omega$ is
\[
Q^*(\omega)=\frac12+\Bigl(\Qmaj-\frac12\Bigr)\min\Bigl\{1,\frac{2(\Gs-\eta)}{\omega Mb}\Bigr\}.
\]
The damped majority rule of strength
$\lambda^*=\min\{1,2(\Gs-\eta)/(\omega Mb)\}$ with majority sharing attains it.
At $\omega=0$, the optimum is $Q^*(0)=\Qmaj$, attained with $\lambda^*=1$.
\end{theorem}

\begin{proof}
At $\omega=0$, majority sharing supports full majority allocation because
$\eta\le\Gs$; Lemma~\ref{lem:influence} gives its optimality. We now take
$\omega>0$ and combine the three lemmas. A rule that supports honest evaluation has
$H(a)\le2(G(v)-\eta)/(\omega M)\le2(\Gs-\eta)/(\omega M)$ by
Lemmas~\ref{lem:ic} and~\ref{lem:rent}, and Lemma~\ref{lem:influence} then gives
\[
Q(a)-\frac12\le\Bigl(\Qmaj-\frac12\Bigr)\min\Bigl\{1,\frac{H(a)}b\Bigr\}
\le\Bigl(\Qmaj-\frac12\Bigr)\lambda^*.
\]
For attainment we take the damped majority rule of strength $\lambda$ with
majority sharing. It has $H=\lambda b$ and $G=\Gs$, so by Lemma~\ref{lem:ic} it
supports honest evaluation exactly when $\lambda\le\lambda^*$, and at
$\lambda=\lambda^*$ its quality is
\[
\frac12+\lambda^*\Bigl(\Qmaj-\frac12\Bigr)=Q^*(\omega).\qedhere
\]
\end{proof}

With free evaluation and $\omega>0$, the frontier reads $Q^*(\omega)=\frac12+(\Qmaj-\frac12)\min\{1,K/\omega\}$
with the ownership threshold
\begin{equation}\label{eq:K}
K=\frac{2\Gs}{Mb}=\frac VM\cdot\frac{d(\Qmaj-p)}{npq\,b}.
\end{equation}
For a positive ownership bound the designer sets the damping at
$\lambda^*=\min\{1,K/\omega\}$; nothing else about the rule depends on
$\omega$. In the running example $\Gs=V/24$ and $K=2/9$, and for $\omega=1$ the
frontier is $83/144$. A positive cost lowers the threshold to
$K(1-\eta/\Gs)$, and at $\eta=\Gs$ only the even split remains when $\omega>0$.

\section{Yuma}
\label{sec:yuma}

\subsection{The rule}

With two miners and equal stake, a validator's weights are a pair
$(1-x_i,x_i)$, where $x_i\in[0,1]$ is its weight on miner $1$. For each miner,
Yuma takes the median of the validators' weights on that miner and clips every
weight above the median down to the median. Miner $1$ receives the share of $M$
equal to its total clipped weight divided by the total clipped weight on both
miners. Validators are paid through bonds: validator $i$'s current share of a
miner is its clipped weight divided by that miner's total clipped weight.
If this total is zero, all current shares in that miner are zero. Each round
we form the rate-$\alpha$ average of old bonds and current shares, with
$\alpha\in(0,1]$, and normalize each positive bond column to sum to one; a zero
column remains zero. Using these updated bonds, validator $i$ receives
$V\bigl(B_{i1}a+B_{i0}(1-a)\bigr)$, where $a$ is miner $1$'s share. Subtensor's
default is $\alpha=1/10$; at $\alpha=1$ bonds equal current shares, and we say
that bond memory is switched off. We call \emph{$\ell$-honest}, for
$\ell\in[0,1/2)$, the strategy of evaluating and submitting weight $1-\ell$ on
the miner one's signal names and $\ell$ on the other.

Table~\ref{tab:round} replays the round of the introduction in the running
example, with Alice, Bob and Carol as validators A, B and C.
Validator C's evaluation favors miner $0$, and the other two validators split.
If C puts all its weight on miner $0$, miner $0$ receives the whole miner
budget. If C moves weight $1/5$ to miner $1$, that weight becomes the median on
miner $1$, the clipping of A's weight stops at $1/5$, and miner $1$ receives a
fifth of the budget. C's own pay stays at $V/2$, and $V/10$ of the validator
budget moves from B to A. A validator that owns part of miner $1$ gains from
this hedge at no cost in this profile; Theorem~\ref{thm:hedging} weighs that
gain against what the hedge costs in the other profiles.

\begin{table}[t]
\caption{One round of Yuma in the running example with bond memory switched
off. Validator A favors miner $1$ and B favors miner $0$; C's evaluation favors
miner $0$, and C hedges by putting weight $1/5$ on miner $1$. Pairs list
(miner $1$, miner $0$). The last row is C's all-or-nothing report.}
\label{tab:round}
\small
\begin{tabular}{@{}lccc@{}}
\toprule
& A & B & C \\
\midrule
weights & $(1,0)$ & $(0,1)$ & $(\frac15,\frac45)$ \\
medians & \multicolumn{3}{c}{$\frac15$ on miner $1$, $\frac45$ on miner $0$} \\
clipped weights & $(\frac15,0)$ & $(0,\frac45)$ & $(\frac15,\frac45)$ \\
current shares & $(\frac12,0)$ & $(0,\frac12)$ & $(\frac12,\frac12)$ \\
pay, miner $1$ gets $\frac15$ & $\frac{V}{10}$ & $\frac{2V}{5}$ & $\frac V2$ \\
\addlinespace
pay if C reports $(0,1)$ & $0$ & $\frac V2$ & $\frac V2$ \\
\bottomrule
\end{tabular}
\end{table}

\begin{proposition}[All-or-nothing weights]\label{prop:binary}
With bond memory switched off, $0$-honest reports make Yuma pay miners by the
majority rule and validators by majority sharing. Hence $0$-honesty is an
equilibrium against all reports in $\{0,1\}$ if and only if
$\eta+\frac12\omega Mb\le\Gs$; with free evaluation, if and only if
$\omega\le K$.
\end{proposition}

\begin{proof}
We take a profile in which $k>m$ validators put all weight on miner $1$. The
median weight is $1$ on miner $1$ and $0$ on miner $0$. Clipping removes every
weight the minority placed on miner $1$ and every weight the majority placed on
miner $0$, so miner $1$ receives $M$ and each majority validator holds a
current share $1/k$ of miner $1$, while the minority holds nothing of it.
The zero column for miner $0$ contributes no payment. The payments are $V/k$ to each majority validator and $0$ to the others, and
unanimity pays $V/n$ each. Lemma~\ref{lem:ic} with $H=b$ and $G=\Gs$ turns the
equilibrium condition into
\[
\eta+\tfrac12\,\omega Mb\le\Gs .\qedhere
\]
\end{proof}

A symmetric honest equilibrium of Yuma induces a rule in $\Rset$, so
Theorem~\ref{thm:frontier} bounds it.

\begin{lemma}[Yuma stays inside the frontier]\label{lem:inherit}
Let bond memory be switched off and let $\ell$-honesty be an equilibrium for
some $\ell\in[0,1/2)$ at cost $\eta$ and ownership bound $\omega$. The map from
the number of validators whose signal names miner $1$ to Yuma's allocation and
payments is a rule in $\Rset$ that supports honest evaluation. In particular,
the quality of $\ell$-honesty is at most $Q^*(\omega)$.
\end{lemma}

The proof (Appendix~\ref{app:yuma}) checks monotonicity of the induced
allocation and notes that every deviation of the induced rule is a deviation in
Yuma.

\begin{theorem}[Hedging]\label{thm:hedging}
Let bond memory be switched off and evaluation be free. For
$k\in\{0,\dots,2m\}$ let
$w_k=\binom{2m}{k}\bigl(p^{2m-k+1}q^k+q^{2m-k+1}p^k\bigr)$, the probability
that $k$ of the other $2m$ validators report $1$ given that one's own signal
is $0$, and let
\[
S_n(p)=\sum_{k<m}w_k\frac{n-k-1}{(n-k)^2}-\sum_{k>m}\frac{w_k}{k}.
\]
Then $0$-honesty is an equilibrium against all weights in $[0,1]$ if and only
if $\omega\le\KY$, where $\KY=VS_n(p)/(Mb)$. Moreover $\KY<K$, and for $n=3$,
\[
\KY=\frac VM\cdot\frac{4-21pq}{36pq}.
\]
\end{theorem}

\begin{table}[t]
\caption{Pay of a validator with signal $0$ that puts weight $x$ on miner $1$
while the others are $0$-honest; $k$ of the other $2m$ report $1$.}
\label{tab:cases}
\small
\begin{tabular}{@{}llcc@{}}
\toprule
Others & Validator's side & Miner $1$'s share & Payment \\
\midrule
$k<m$ & majority & $0$ & $V(1-x)/(n-k-x)$ \\
$k=m$ & pivotal & $x$ & $V/(m+1)$ \\
$k>m$ & minority & $1$ & $Vx/(k+x)$ \\
\bottomrule
\end{tabular}
\end{table}

The three rows of Table~\ref{tab:cases} are the hedging account of the
introduction: a small weight on miner $1$ costs a little when the validator
sides with the majority, earns a share of miner $1$'s bonds when it sides with
the minority, and moves the allocation one-for-one when it is pivotal.

\begin{table}[t]
\caption{Raw threshold ratios $\KY/K$ from the closed forms, with $V=M$,
free evaluation and bond memory off. Negative values mean that a small
hedge pays at zero ownership. A dagger marks $K,\KY\ge1$: both then support
fully decisive reports throughout the feasible ownership range $[0,1]$.}
\label{tab:thresholds}
\centering
\small
\input{tables/thresholds}
\end{table}

Table~\ref{tab:thresholds} extends the comparison beyond the running example.
At $p=3/4$, the ratio grows with the displayed committee sizes, and already
at $n=15$ both raw thresholds exceed one. At $p=3/5$, $\KY$ is negative
for $n=3,5,\ldots,13$, but positive at $n=15$ ($\KY\approx0.00998$).
Thus zero-ownership hedging is a small-committee result at this accuracy,
not a statement for every $n$. The table evaluates the formulas exactly
before rounding; it assumes equal stake and conditionally independent signals.

Softer weights let Yuma keep some information above $\KY$.

\begin{proposition}[Softer weights]\label{prop:soft}
Let $n=3$, $V=M$, $p=3/4$, bond memory switched off and evaluation free. For
$\ell\in[0,1/2)$, $\ell$-honesty is an equilibrium against all weights in
$[0,1]$ if and only if $\omega\le\tau(\ell)$, where
\[
\tau(\ell)=\frac{2(28\ell^2+56\ell+1)}{3(28\ell^3+56\ell^2+127\ell+72)} .
\]
The tolerance $\tau$ increases from $\tau(0)=1/108$ to $\tau(1/2)=8/51$, and the
quality of $\ell$-honesty,
$\QY(\ell)=\frac12+\frac12\bigl(\frac12-\ell\bigr)\bigl(\frac{13}{16}+\frac{9}{16(1+\ell)}\bigr)$,
decreases from $27/32$ to $1/2$.
\end{proposition}

At each $\omega\in(1/108,8/51)$ the best $\ell$-honest equilibrium therefore has
quality $\QY(\tau^{-1}(\omega))<27/32=Q^*(\omega)$. Softer weights damp Yuma's
allocation, but they spread a report's influence over profiles in which no
single report decides the majority, where influence buys less quality
(Lemma~\ref{lem:influence}). The damped majority rule keeps all influence in the
profile where the others split evenly. Figure~\ref{fig:influence} shows the
difference with seven validators: at equal influence, a quarter of Yuma's
influence sits in profiles where the other validators already form a majority,
and its quality falls short of the damped majority rule's by $0.034$.

\begin{figure*}[t]
\centering
\includegraphics[width=\textwidth]{figures/influence.pdf}
\caption{Allocation and influence with seven validators and $p=3/4$. (a) Miner
$1$'s share against the number of reports for miner $1$ under the majority
rule, under Yuma with $\frac15$-honest weights (bond memory off), and under the
damped majority rule of strength $\lambda=0.62$, which has the same influence
$H=0.082$ as Yuma. (b) The contribution of each count $j$ of other reports for
miner $1$ to the influence~\eqref{eq:influence}. The damped majority rule places
all of its influence where the other six validators split evenly and reaches
quality $0.766$; Yuma places a quarter of it where no single report decides the
majority and reaches $0.731$.}
\label{fig:influence}
\end{figure*}

\subsection{Default bond memory}

With $\alpha<1$ a validator's current weights also buy future bonds, and the
game becomes an infinitely repeated one.

\begin{theorem}[Default parameters]\label{thm:default}
Let $n=3$, $p=3/4$, $V=M=1/2$, $\alpha=1/10$, discount factor $\delta=99/100$,
the state $\theta$ drawn independently in each round, and every initial bond
equal to $1/3$. In the repeated game:
\begin{enumerate}
\item For every $\ell\in[0,2/5]$, $\ell$-honesty is not a Nash equilibrium when
all ownership is zero.
\item $\tfrac{49}{100}$-honesty is a Nash equilibrium whenever
$|\beta_i|\le1/2000$ for all $i$ and $0\le\eta\le1/400000$. Its quality is
$241237/476800$.
\item Not evaluating and submitting weight $1/2$ on each miner in every round
is a Nash equilibrium for every $\eta\ge0$ and every ownership profile with
$|\beta_i|\le1$.
\end{enumerate}
\end{theorem}

The term $(1-\alpha)D(I-\frac12)$, with bond difference $D=B_{i1}-B_{i0}$
and miner $1$'s share $I$, explains part~(i). A validator whose bonds
lean toward miner $1$ by $D$ is paid like an owner of a share
$(1-\alpha)DV/M$ of miner $1$, so its old bonds act as ownership of the miners
it has backed. Along honest play the lean reaches
$d_*=2(1-2\ell)/(3(2-\ell))$, which is $1/3$ for all-or-nothing weights and
$4/453$ at $\ell=49/100$. Decisive honest weights thus build the ownership that
Theorem~\ref{thm:hedging} shows they cannot withstand, while nearly even
weights keep the lean small.

\section{Unequal stake and persistence}
\label{sec:extensions}

\subsection{Unequal stake}

Let validator $i$ hold positive stake share $\varphi_i>0$, with $\sum_i\varphi_i=1$. The
allocation now depends on the reports through the stake that reports $1$,
$a=f(\sum_i\varphi_ir_i)$, with $f$ nondecreasing and $f(1-z)=1-f(z)$. Payments
may depend on identities: $v_i(r)\ge0$, $\sum_iv_i(r)=V$ for every report
profile $r$, and $v_i(r)=v_i(\mathbf 1-r)$. Validator $i$ has its own premium
$G_i(v)=\E[v_i(s)]-\E[v_i(1,s_{-i})]$ and coordinate influence $H_i(a)$
as in Section~\ref{sec:model}. The argument of Lemma~\ref{lem:ic} gives the constraint
$\eta+\frac12\omega MH_i(a)\le G_i(v)$ for each $i$. Which premium vectors can
a fixed budget deliver?

\begin{theorem}[Premium capacity]\label{thm:capacity}
For $j=m+2,\dots,n$ let
$R_j=\frac{Vd}{2}\bigl(p^{j-1}q^{n-j}-q^{j-1}p^{n-j}\bigr)$, and for
$k=0,\dots,n$ let
\[
C_k=\sum_{j=m+2}^{n}\Bigl[\binom nj-\binom{n-k}{j}\Bigr]R_j .
\]
A vector $(g_1,\dots,g_n)\ge0$ satisfies $G_i(v)\ge g_i$ for all $i$ under
some payment rule if and only if $\sum_{i\in T}g_i\le C_{|T|}$ for every set
$T$ of validators. The total capacity is $C_n=n\Gs$.
\end{theorem}

Each profile in which a majority of $j\ge m+2$ validators agrees creates
premium $R_j$ that only its members can receive; a group of size $k$ can draw
on every such majority it meets. Combined with a decomposition of stake-based
allocations into threshold rules, Theorem~\ref{thm:capacity} turns the frontier
under unequal stake into a linear program with an explicit payment
construction. With three validators, one holding at least half the stake and
the other two equal, let $0\le\eta\le Vd^2/6$. For $\omega>0$ the closed form is
$Q^*=\frac12+\frac d2(1+pq)\min\{1,(Vd^2-6\eta)/(\omega M(\frac12+3pq))\}$;
at $\omega=0$ it is $\frac12+\frac d2(1+pq)$, and above the cost bound no rule
is feasible (Appendix~\ref{app:extensions}). In the running example the largest validator lowers the achievable quality from
$27/32$ to $51/64$, and the ownership threshold becomes $4/17$.

Yuma divides the validator budget among the winning side in proportion to
stake. In the running example with stake shares $(3/5,1/5,1/5)$ and all-or-nothing
weights, the largest validator's premium under this rule is $-V/20$: it decides
the majority alone, keeps a larger share of the budget when fewer validators
side with it, and so prefers a blind report even without ownership.
Theorem~\ref{thm:capacity} shows that the same budget can give this validator a
premium up to $C_1=Vd^2/2=V/8$.

\subsection{Persistence}

Now the state is a stationary binary symmetric Markov chain, with uniform
initial state and transition probability
$\Pr[\theta_{t+1}=\theta_t\mid\theta_0,\ldots,\theta_t]=(1+\rho)/2$
for $\rho\in[0,1)$. Signals follow Section~\ref{sec:model}. In each round validators may evaluate again
at cost $\eta>0$, reports become public after the round, and the rule is a
fixed element of $\Rset$ applied to the current round's reports. Validators
discount the future and may deviate using their entire private and public
history.

\begin{theorem}[Frontier with persistence]\label{thm:memory}
Let $\Lambda=(p/q)^n$ and let $\Pi\in[1/2,1)$ solve
$\Pi=\frac{1-\rho}2+\rho\frac{\Lambda\Pi}{1+(\Lambda-1)\Pi}$.
If $\Pi\ge p$, no rule of this kind supports honest evaluation. If $\Pi<p$, let
$G_\rho=2(p-\Pi)\Gs/d$; honest evaluation can be supported if and only if
$\eta\le G_\rho$. On this feasible domain the largest quality is
$Q^*_\rho(0)=\Qmaj$ at zero ownership; for $\omega>0$ it is
\[
Q^*_\rho(\omega)=\frac12+\Bigl(\Qmaj-\frac12\Bigr)\min\Bigl\{1,\frac{G_\rho-\eta}{\omega Mb\,(p-d\Pi)}\Bigr\},
\]
attained by repeating the damped majority rule with majority sharing. The
condition $\Pi\ge p$ holds if and only if
$\rho\ge\rho_c=d\,\frac{p^{n+1}+q^{n+1}}{p^{n+1}-q^{n+1}}$, and $\rho_c$
decreases to $d$ as $n$ grows.
\end{theorem}

The number $\Pi$ is the most confident belief a validator can hold before
evaluating, reached after long runs of unanimous reports. A validator whose
belief is $\pi$ gains $2(p-\pi)G(v)/d$ from evaluating, which vanishes as $\pi$
approaches $p$: at that point a fresh signal can no longer overturn what the
history already says. The constraint binds at the most confident history. In
the running example $\rho_c=41/80$, so fresh evaluation stops paying once the
better miner stays better with probability $(1+\rho_c)/2=121/160$.

\section{The deployed rule}
\label{sec:replica}

The analysis of Section~\ref{sec:yuma} uses exact arithmetic. Subtensor
computes Yuma in fixed point, stores weights and bonds as 16-bit integers, and
normalizes them at several steps. To check that these details do not move the
thresholds, we ported the production sparse epoch function of
Subtensor~\cite{subtensor2026} (the legacy path, the path with Yuma3 bonds, and
the path with adaptive bond rates) to Python with integer arithmetic that
follows the fixed-point semantics of the source. As a reference we compiled the
original Rust functions with their arithmetic unchanged, reading chain state
through a small adapter. The original
mathematical test suite passes (81 tests), as do the six official Yuma3 test
sequences (36 epochs). The Python port agrees with the Rust reference bit for
bit on 516 states: the official test inputs, 80 boundary states and 400 random
states. Table~\ref{tab:replica} compares the closed forms of
Section~\ref{sec:yuma} with the implementation. For each $\ell$, we sampled
about $1040$ fixed-sum 16-bit reports, including nearby boundary reports,
for both signals and both ownership signs. The reported cap deters the
sampled deviations; it is not a certificate for all 16-bit reports.
The sampled caps differ from the closed forms by less than
$2.5\times10^{-5}$; Appendix~\ref{app:replica} gives all five values and
the grid construction.

\paragraph{One validator, multiple miners.} We next retain the full
fixed-point rule with bond memory off and vary the number of validators
and miners. Each validator observes independent Gaussian scores: the
best miner has mean $1.3$, the others mean $1$, and the noise standard
deviation is $0.3$. Signal weights are the softmax of these scores at
temperature $0.3$. Only validator 0 deviates, replacing its weight vector
$w$ by $(1-h)w+h(1/J,\ldots,1/J)$, where $J$ is the number of miners.
It owns no miner and still evaluates, so the gain is entirely in its
validator payment.

Table~\ref{tab:single-hedge} reports the preselected $h=0.1$ comparisons,
using 2048 paired draws per environment. The remaining validators split
the stake outside $s_0$ equally. The complete grid, including intermediate
stakes and $h=0.25$, is in the replication data. All 32 softmax comparisons have
positive payment gains and negative quality changes, with their pointwise
95\% intervals excluding zero. Measured against its pay without the hedge,
an equal-stake validator ($s_0=1/n$) gains 1.8\% to 4.8\% with 11 validators
and 4.5\% with 31. A 40\% stake holder gains 0.3\% to 0.6\% of its pay, and
its hedge lowers quality by 1.1 to 1.6 percentage points. These are deviations from the specified signal-weight
policy, not an equilibrium classification or a test of patient behavior
under default bond memory. All-or-nothing controls include negative gains
(Appendix~\ref{app:replica}).

\begin{table}[t]
\caption{A single zero-owner validator hedges with $h=0.1$.
$\Delta v/V$ is its payment gain as a fraction of the validator budget;
$\Delta Q$ is the change in the share paid to the best miner.
Payoff and quality entries are multiplied by $10^3$; brackets are paired, pointwise
95\% Monte Carlo intervals. Bond memory is off.}
\label{tab:single-hedge}
\centering
\small
\input{tables/single_hedge_main}
\end{table}

\section{Conclusion}
\label{sec:conclusion}

A reward rule that sees only reports buys honest evaluation with a premium
that a fixed budget caps, and a validator's ownership prices every unit of
influence its report carries. The best rule spends the premium on agreement
about the verdict and damps the verdict's weight in the allocation once
ownership passes $K$. Yuma reaches this rule only when weights are
all-or-nothing. With continuous weights it also pays for partial agreement,
which is what a hedging validator supplies. A subnet can remove the hedging
channel in the two-miner model by allocating and rewarding on the miner
each validator weights most, sharing the validator budget among the majority,
and damping the allocation by $\min\{1,K/\omega\}$.

\bibliographystyle{ACM-Reference-Format}
\bibliography{references}

\appendix

\section{Proofs for Section~\ref{sec:frontier}}
\label{app:frontier}

Write $h_j=\frac12\binom{n-1}{j}\bigl(p^jq^{n-1-j}+q^jp^{n-1-j}\bigr)$ for the
probability that $j$ of the other $n-1$ validators report $1$ when all evaluate
honestly and $\theta$ is uniform, so that $\E[v_{1,1+J}]=\sum_jh_jv_{1,1+j}$.

\begin{proof}[Proof of Lemma~\ref{lem:ic}]
We check the two kinds of deviation in turn, starting with a validator who
evaluated and holds signal $s$. The symmetry
$v_{0,k}=v_{1,n-k}$ makes the expected payment of reporting $s$ exceed that of
reporting $1-s$ by $2G(v)$ for either signal. The increments
$a_{k+1}-a_k$ are symmetric under $k\mapsto n-1-k$, so switching the report
changes the expected share of miner $1$ by $H(a)$ conditional on either signal.
Against the worst sign of ownership, honest reporting after evaluation is
optimal if and only if
\begin{equation}\label{eq:after}
2G(v)-\omega MH(a)\ge0 .
\end{equation}
A validator who skips the evaluation earns the same expected payment from
either fixed report. Under honest evaluation miner $1$'s expected share is
$1/2$, and a fixed report for miner $1$ raises it by $H(a)/2$. Skipping also
saves $\eta$, so evaluation beats every blind report if and only if
$G(v)-\frac12\omega MH(a)-\eta\ge0$. A randomized blind report is a mixture of
the two fixed reports, randomized reports after evaluation are mixtures of pure
reports for each signal, and randomizing the evaluation decision cannot beat
the better pure decision. Because $\eta\ge0$, the last condition already
implies~\eqref{eq:after}:
\[
2G(v)-\omega MH(a)\ge2\eta\ge0 .\qedhere
\]
\end{proof}

\begin{proof}[Proof of Lemma~\ref{lem:rent}]
By the budget~\eqref{eq:budget} and the equal treatment of validators, the
premium is $V/n$ minus the blind payment $\sum_jh_jv_{1,1+j}$, so we minimize
the latter. Write $t_k=v_{1,k}$. The symmetries give
$kt_k+(n-k)t_{n-k}=V$. For $1\le k\le m$ we eliminate $t_{n-k}$ from each pair;
using $h_j=h_{n-1-j}$, the coefficient of $t_k$ in the blind payment is
\[
h_{k-1}-\frac{k}{n-k}\,h_k=\frac12\binom{n-1}{k-1}(p-q)\bigl(q^{k-1}p^{n-k-1}-p^{k-1}q^{n-k-1}\bigr)>0,
\]
since $p>q$ and $n-k-1>k-1$. We therefore set $t_k=0$ and $t_{n-k}=V/(n-k)$ for
$1\le k\le m$; unanimity fixes $t_n=V/n$. This is majority sharing. Its blind
payment follows from $\frac1k\binom{n-1}{k-1}=\frac1n\binom nk$:
\[
\sum_jh_j\,v_{1,1+j}=\frac{V}{2n}\Bigl(\frac{\Qmaj}{p}+\frac{1-\Qmaj}{q}\Bigr),
\]
and therefore
\[
\frac Vn-\frac{V}{2n}\Bigl(\frac{\Qmaj}{p}+\frac{1-\Qmaj}{q}\Bigr)=\frac{Vd(\Qmaj-p)}{2npq}.
\qedhere
\]
\end{proof}

\begin{proof}[Proof of Lemma~\ref{lem:influence}]
The first bound holds profile by profile: given $k$ reports for miner $1$, the
posterior that miner $1$ is better exceeds $1/2$ exactly when $k>n/2$, so the
majority rule pays the likelier miner in every profile. For the second bound
we decompose the allocation into threshold rules. Let
$\Delta_j=a_{j+1}-a_j\ge0$; symmetry gives $\Delta_j=\Delta_{2m-j}$ and
$\sum_j\Delta_j\le1$. For $j<m$ let
$f_j(k)=\frac12(\one\{k>j\}+\one\{k>2m-j\})$, and let $f_m(k)=\one\{k>m\}$.
Then
\[
a=\Bigl(1-\sum_j\Delta_j\Bigr)\frac12+\sum_{j<m}2\Delta_j\,f_j+\Delta_m f_m ,
\]
and $Q$ and $H$ are affine in $a$, so we may compare the thresholds one at a
time. For an accuracy $t\in[1/2,p]$ let $Q_j(t)=\E f_j(\Bin(n,t))$ and let
$H_j(t)$ be the influence of $f_j$ at accuracy $t$:
\begin{gather*}
H_j(t)=\tfrac12\tbinom{2m}{j}\bigl[t^j(1-t)^{2m-j}+t^{2m-j}(1-t)^j\bigr]\quad(j<m),\\
H_m(t)=\tbinom{2m}{m}\bigl[t(1-t)\bigr]^m .
\end{gather*}
Differentiating the binomial sums gives
\[
Q_j'(t)=nH_j(t),\qquad Q_j(1/2)=1/2,
\]
and the ratio of the two influences is
\[
\frac{H_j(t)}{H_m(t)}=\frac{\tbinom{2m}{j}}{\tbinom{2m}{m}}\,\cosh\Bigl((m-j)\log\tfrac{t}{1-t}\Bigr),
\]
which increases for $t\ge1/2$. Hence
\[
\frac{H_j(t)}{H_j(p)}\le\frac{H_m(t)}{H_m(p)}\qquad\text{for }t\in[1/2,p],
\]
and since $H_m(p)=b$ and $Q_m(p)=\Qmaj$, integrating gives
\begin{align*}
\frac{Q_j(p)-1/2}{H_j(p)}&=n\int_{1/2}^{p}\frac{H_j(t)}{H_j(p)}\,dt\\
&\le n\int_{1/2}^{p}\frac{H_m(t)}{H_m(p)}\,dt=\frac{\Qmaj-1/2}{b}.
\end{align*}
With mixture weights $c_j=2\Delta_j$ for $j<m$ and $c_m=\Delta_m$, we conclude
\begin{align*}
Q(a)-\frac12&=\sum_{j\le m}c_j\Bigl(Q_j(p)-\frac12\Bigr)
\le\frac{\Qmaj-1/2}{b}\sum_{j\le m}c_jH_j(p)\\
&=\frac{\Qmaj-1/2}{b}\,H(a).\qedhere
\end{align*}
\end{proof}

\section{Proofs for Section~\ref{sec:yuma}}
\label{app:yuma}

\begin{proof}[Proof of Lemma~\ref{lem:inherit}]
We first check that the induced allocation is monotone. Under $\ell$-honesty,
when $k>m$ validators name miner $1$ the median weight on
miner $1$ is $1-\ell$ and on miner $0$ is $\ell$. Clipping leaves miner $1$ with
total weight $k(1-\ell)+(n-k)\ell$ and miner $0$ with $n\ell$, so
\[
a_k=\frac{k(1-\ell)+(n-k)\ell}{k+2\ell(n-k)},\qquad
\frac{\partial a_k}{\partial k}=\frac{(1-2\ell)\,n\ell}{\bigl(k+2\ell(n-k)\bigr)^2}\ge0 ,
\]
and $a_{n-k}=1-a_k$ by symmetry, so $a_m\le\frac12\le a_{m+1}$. Payments sum to
$V$ in every profile, are nonnegative, and treat validators and miners alike.
The induced rule is therefore in $\Rset$. Each deviation of the induced rule
(evaluating and reporting against one's signal, or reporting a fixed miner
without evaluating) corresponds to submitting the weights of an
$\ell$-honest validator with the other signal, which is available in Yuma.
Hence the induced rule supports honest evaluation, and
Theorem~\ref{thm:frontier} gives $\QY(\ell)\le Q^*(\omega)$.
\end{proof}

\begin{proof}[Proof of Theorem~\ref{thm:hedging}]
Let the validator have signal $0$ and weight $x$ on miner $1$. When $k<m$
others report $1$, the medians are $0$ on miner $1$ and $1$ on miner $0$, so
the validator's weight on miner $1$ is clipped away, miner $0$ receives $M$,
and the validator holds $(1-x)/(n-k-x)$ of miner $0$. When $k>m$, it holds
$x/(k+x)$ of miner $1$, which receives $M$. When $k=m$, its weight $x$ is the
median on miner $1$, the other $m$ supporters of miner $1$ are clipped to $x$,
miner $1$ receives $xM$, and the validator holds $1/(m+1)$ of each miner
with positive clipped weight. At $x=0$ or $1$, the zero column contributes
nothing; the payment is still $V/(m+1)$. This gives Table~\ref{tab:cases}. The validator's signal-$0$ gain of weight $x$ over
weight $0$ is
\[
F_\beta(x)=Vx\Bigl[\sum_{k>m}\frac{w_k}{k+x}-\sum_{k<m}\frac{w_k(n-k-1)}{(n-k)(n-k-x)}\Bigr]+\beta Mw_mx,
\]
and $w_m=b$. The maps $x\mapsto x/(k+x)$ and $x\mapsto(1-x)/(n-k-x)$ are
strictly concave on $[0,1]$, so $F_\beta$ is strictly concave, and $F_\beta(0)=0$.
Weight $0$ is optimal against every $|\beta|\le\omega$ if and only if
$F_\omega'(0)=-VS_n(p)+\omega Mb\le0$, that is, $\omega\le\KY$. The same holds
for signal $1$ by symmetry. With free evaluation, a validator who does not
evaluate and submits weight $x$ earns the average of the payoffs that weight
$x$ earns under the two signals, and each of these is at most the honest payoff
for that signal; so skipping the evaluation does not pay either. To compare
$\KY$ with $K$, note that $F_0(1)$ is the payment loss from reporting against
one's signal, which is
$-2\Gs$ under majority sharing (proof of Lemma~\ref{lem:ic}), and strict
concavity gives $F_0'(0)>F_0(1)-F_0(0)=-2\Gs$, so $VS_n(p)<2\Gs$ and $\KY<K$.
For $n=3$ we have $w_0=p^3+q^3=1-3pq$, $w_2=pq$ and $b=2pq$, so
\[
S_3(p)=\frac{2w_0}{9}-\frac{w_2}{2}=\frac{4-21pq}{18},
\qquad
\KY=\frac{VS_3(p)}{Mb}=\frac VM\cdot\frac{4-21pq}{36pq}.
\qedhere
\]
\end{proof}

\begin{proof}[Proof of Proposition~\ref{prop:soft}]
We fix a validator with signal $0$, write $t=pq$ and $u=1-3t$, and let the
other two validators be $\ell$-honest. Given the validator's signal, the other
two put weight $\ell$, one $\ell$ and one $1-\ell$, or $1-\ell$ on miner $1$
with probabilities $u$, $2t$ and $t$. Let $v(x)$ be the validator's expected
payment divided by $V$ and $i(x)$ miner $1$'s expected share when it puts weight
$x$ on miner $1$. Computing the medians and the clipping as in
Table~\ref{tab:round} gives, for $x\le\ell$,
\begin{align*}
v_L&=u\tfrac{1-\ell+x}{3-\ell+x}+t\tfrac{2-\ell+3x}{2+\ell+x},\\
i_L&=u\tfrac{x+2\ell}{3-\ell+x}+t\tfrac{2+2\ell+3x}{2+\ell+x};
\end{align*}
for $\ell\le x\le1-\ell$,
\begin{align*}
v_C&=u\tfrac{1-x+\ell}{3-x+\ell}+t\tfrac{\ell+x}{2+\ell+x}+\tfrac{t}{1+\ell},\\
i_C&=\tfrac{3u\ell}{3-x+\ell}+t\bigl(1-\tfrac{3\ell}{2+\ell+x}\bigr)+\tfrac{t(\ell+2x)}{1+\ell};
\end{align*}
and for $x\ge1-\ell$,
\begin{align*}
v_R&=\tfrac{u(1-x+\ell)+2t(2-\ell-x)}{3-x+\ell}+t\tfrac{2-\ell-x}{4-\ell-x},\\
i_R&=\tfrac{3u\ell+2t(2-\ell)}{3-x+\ell}+\tfrac{3t(1-\ell)}{4-\ell-x}.
\end{align*}
The pieces agree at the interfaces, and $v_C(\ell)=1/3$. With $V=M$ and
ownership $\beta$, the gain of weight $x$ over weight $\ell$ is
$V\{v(x)-v(\ell)+\beta[i(x)-i(\ell)]\}$. Since $i$ increases in $x$, the worst
ownership is $+\omega$ for moves to the right and $-\omega$ for moves to the
left.

\emph{Necessary bound.} On the middle interval, with $\beta=\omega$, the derivative of the gain is
\[
\frac{u(-2+3\omega\ell)}{(3-x+\ell)^2}+\frac{t(2+3\omega\ell)}{(2+\ell+x)^2}+\frac{2\omega t}{1+\ell}.
\]
At $x=\ell$ its coefficient of $\omega$ is positive. A necessary condition
for equilibrium is therefore $\omega\le C_C(\ell)$, where
\[
C_C(\ell)=\frac{2u/9-t/[2(1+\ell)^2]}{u\ell/3+3\ell t/[4(1+\ell)^2]+2t/(1+\ell)} .
\]
For $p=3/4$, this bound is $C_C(\ell)=\tau(\ell)\le8/51<1$.
\emph{Sufficiency.} We henceforth take $0\le\omega\le C_C(\ell)$.
The middle derivative then has a negative first coefficient and a positive
second coefficient, so it decreases in $x$ and remains nonpositive.
On the left interval, with $\beta=-\omega$, the derivative is
$A_L/(3-\ell+x)^2+B_L/(2+\ell+x)^2$ with $A_L=u[2-3\omega(1-\ell)]$ and
$B_L=t[4(1+\ell)-\omega(4+\ell)]\ge3t\ell$. If $A_L\ge0$ it is positive. If
$A_L<0$ we multiply by $(2+\ell+x)^2$; the negative term then carries the factor
$[(2+\ell+x)/(3-\ell+x)]^2$, which increases in $x$, so nonnegativity at
$x=\ell$ controls the interval and gives $\omega\le C_L(\ell)$ with
\[
C_L(\ell)=\frac{2u/9+t/(1+\ell)}{u(1-\ell)/3+t(4+\ell)/[4(1+\ell)^2]} .
\]
On the right interval, with $\beta=\omega$, the derivative is
$C_R/(3-x+\ell)^2+D_R/(4-\ell-x)^2$ with
$C_R=u(-2+3\omega\ell)+2t[-(1+2\ell)+\omega(2-\ell)]$ and
$D_R=t[-2+3\omega(1-\ell)]$. The coefficient $C_R$ increases in $\omega$ and
equals $-(1-4t)(2-3\ell)-3t\ell\le0$ at $\omega=1$. The derivative jumps down at
$x=1-\ell$, so it is nonpositive just right of the interface whenever the
middle condition holds; if $D_R\le0$ it stays nonpositive, and if $D_R>0$ the
factor $[(3-x+\ell)/(4-\ell-x)]^2$ multiplying the positive term after scaling
by $(3-x+\ell)^2$ decreases in $x$, which again keeps the sign. The right
interval therefore adds no condition. With free evaluation, skipping the
evaluation does not pay once the reporting conditions hold, by the averaging
argument in the proof of Theorem~\ref{thm:hedging}. At $p=3/4$ we have
$C_L>C_C$ on $[0,1/2)$ and $C_C=\tau$, and the allocation formula of
Lemma~\ref{lem:inherit} gives
\[
\QY(\ell)=\frac12+d\Bigl(\frac12-\ell\Bigr)\Bigl(1-t+\frac{3t}{1+\ell}\Bigr).\qedhere
\]
\end{proof}

\begin{proof}[Proof of Theorem~\ref{thm:default}]
For parts~(i) and~(ii), we fix the others at $\ell$-honesty with $\ell>0$;
the zero endpoint is treated separately. We write $r=1-\alpha$ and
$z=\delta r$, let $B_1,B_0$ be the deviator's
bonds, and set $S=B_1+B_0$ and $D=B_1-B_0$. For weight $x$ and a count $k$ of
other reports for miner $1$, let $I(x,k)$ be miner $1$'s share and
$\Delta_j(x,k)$ the deviator's current share of miner $j$. The other validators
put weight at least $\ell>0$ on each miner, so every miner has positive clipped
weight and the bonds evolve as $B_j'=rB_j+\alpha\Delta_j$. With three
equal-stake validators a clipped weight never exceeds the sum of the other two,
so $\Delta_j\le1/2$ and $|D|\le1/2$ under every deviation.

\emph{Reduction.} In units of $V$ the deviator's current pay is
$r\{S/2+D(I-\frac12)\}+\alpha\{I\Delta_1+(1-I)\Delta_0\}$. We add the value
$\delta\kappa S'$ and subtract $\kappa S$ with $\kappa=r/(2(1-z))$, which removes
$S$ and leaves the one-round payoff
\begin{equation}\label{eq:reduced}
R(D,x,k)=rD\Bigl(I-\frac12\Bigr)+\alpha\Bigl\{I\Delta_1+(1-I)\Delta_0+\frac{z(\Delta_1+\Delta_0)}{2(1-z)}\Bigr\}.
\end{equation}

\emph{Potential.} Let $\psi(D,x)$ be the expectation of $R$ given signal $0$
and $\Phi(D)=\frac1{20}(|D|-\frac1{20})_+$. For $\ell=49/100$, all
$D\in[-\frac12,\frac12]$ and all $x\in[0,1]$,
\begin{equation}\label{eq:potential}
\psi(D,x)-\psi(D,\ell)+\delta\,\E\Phi(D')-\Phi(D)\le-\varepsilon|x-\ell|,\qquad \varepsilon=\frac1{1000}.
\end{equation}
For fixed $x$ the left side is convex in $D$ on each of the intervals with
endpoints $-\frac12,-\frac1{20},\frac1{20},\frac12$, so the four endpoints
suffice. On each of the report intervals $[0,\ell]$, $[\ell,1-\ell]$ and
$[1-\ell,1]$ the quantities $I$ and $\Delta_j$ are rational in $x$, and the
three positive parts in $\E\Phi(D')$ split into $27$ affine branches. The
resulting $4\cdot3\cdot27=324$ univariate rational inequalities have positive
denominators and nonnegative Bernstein coefficients on their intervals. The
inequality for signal $1$ follows by symmetry.

\emph{All deviations.} Averaging~\eqref{eq:potential} over the two signals, a
deviator who evaluates satisfies $\E[R-c+\delta\Phi(D')]-\Phi(D)\le C-c$, where
$c=\eta/V$ and $C=\alpha/(3(1-z))$ is the average payoff of $\ell$-honesty. A
deviator who does not evaluate cannot condition $x$ on the signal, and
$\frac12(|x-\ell|+|x-1+\ell|)\ge\frac12-\ell$, so its bound is
$C-\varepsilon(\frac12-\ell)\le C-c$ whenever $c\le\varepsilon(\frac12-\ell)=10^{-5}$.
Multiplying by $\delta^t$ and summing over rounds telescopes the potential.
Since $\Phi(0)=0$ and $\Phi$ is bounded, every deviation earns at most
$(V/3-\eta)/(1-\delta)$, which $\ell$-honesty attains. With ownership
$\beta_i$ the round's payoff gains $\beta_i(M/V)[i(x)-i(\ell)]$. On all nine
branches $0\le\partial I/\partial x\le1$, so this term is at most
$\omega(M/V)|x-\ell|$, and~\eqref{eq:potential} holds with margin
$\varepsilon-\omega M/V\ge1/2000$ when $\omega\le1/2000$ and $V=M$; the margin
covers costs up to $V\cdot\frac1{2000}(\frac12-\ell)=1/400000$. Along honest play
$|D|\le d_*=4/453<1/20$, so the potential vanishes on the equilibrium path. This
proves part~(ii); the quality is $\QY(49/100)$.

\emph{Part (i).} After $T$ rounds in which the deviator's signal is $1$ and the
others' are $0$, an event of probability $(3/32)^T$, the lean is
$D_T=(1-r^T)d_*$ with $d_*=2(1-2\ell)/(3(2-\ell))$. At $D=d_*$ and signal $0$,
the right derivative of $\psi(D,\cdot)$ at $x=\ell$ is a rational function of
$\ell$ whose numerator, after multiplication by $\ell$, has positive Bernstein
coefficients on $[0,\frac25]$, so it is positive for every
$\ell\in(0,\frac25]$. By continuity in $D$ it stays positive at $D_T$ for a finite
$T$, and after that history a slightly larger weight on miner $1$, followed by
$\ell$-honesty, is strictly profitable. For $\ell=0$ the deviation pays in the
first round: with signal $0$ and equal initial bonds, weight $1/100$ on
miner $1$ raises the discounted payoff by $53506799/2288569920$.

\emph{Part (iii).} We fix opponents at $1/2$ and take $\beta_i\ge0$ by
symmetry. Write $Z_j=1/n-B_j$, $Z=Z_1+Z_0$, and $I=1/2+\sigma\xi$, where
$\sigma\in\{-1,1\}$ and $0\le\xi\le\bar\xi=1/(4n-2)$. Clipping gives fresh
shares $1/n$ and $1/n-\chi\xi/(1-2\xi)$, $\chi=4(n-1)/n$, on the favored
and other miner. Thus $Z_j\ge0$ and $Z'=rZ+\alpha\chi\xi/(1-2\xi)$.
We simulate any strategy from the same observations and randomization,
keeping acquisition choices but reflecting reports toward miner $1$.
Fixed opponents make this implementable. Reflection preserves $Z$ and sets
$Z_1=0$; its payoff advantage is
$2V\xi Z_1'\ge0$ if $\sigma=1$, and $2\xi(VZ_0'+\beta_iM)\ge0$ otherwise,
using the original path's updated deficits. For a reflected path the gross
gain is $g=(\beta_iM-\alpha V\chi/2)\xi-VrZ(1/2-\xi)$.
We set $\Psi(Z)=-\zeta Z$, $\zeta=Vr(1/2-\bar\xi)/(1-\delta r)$ and
$\bar\beta=(\alpha V\chi/2+\delta\alpha\zeta\chi)/M$. Then
\begin{align*}
g+\delta\Psi(Z')-\Psi(Z)
&=-VrZ(\bar\xi-\xi)-(\bar\beta-\beta_i)M\xi\\
&\quad-\frac{2\delta\alpha\zeta\chi\xi^2}{1-2\xi}\le0
\end{align*}
for $\beta_i\le\bar\beta$. Since the initial deficits vanish and $\Psi$ is
bounded, discounted summation excludes every such path; acquisition costs
only lower its payoff. At the stated parameters,
$\bar\beta=8218/8175>1$.\qedhere
\end{proof}

\section{Proofs for Section~\ref{sec:extensions}}
\label{app:extensions}

\begin{proof}[Proof of Theorem~\ref{thm:capacity}]
We pair each report profile $r$ with its label flip $\mathbf1-r$ and write $W$
for the set of validators in the majority of $r$, with $|W|=j$. In $G_i(v)$ the
payment $v_i(r)$ carries the coefficient $\Pr(r)-\frac12\Pr(r_{-i})$, which for
$i\in W$ equals
\[
c_j=\frac d4\bigl(p^{j-1}q^{n-j}-q^{j-1}p^{n-j}\bigr),
\]
zero for $j=m+1$ and positive for $j\ge m+2$, and which is negative for
$i\notin W$. Moving a minority member's payment to the majority therefore
weakly raises every premium, so each set $W$ with $|W|=j\ge m+2$ contributes a
resource $2Vc_j=R_j$ that only its members can receive. Majority sharing
splits these resources symmetrically, so $C_n=n\Gs$ by Lemma~\ref{lem:rent}.
A set $T$ of $k$
validators can draw only on the sets $W$ that meet it, and there are
$\binom nj-\binom{n-k}{j}$ of them of each size $j$; this proves necessity. For
sufficiency we fix weights $z\ge0$. The largest $z$-weighted premium is
$\sum_WR_{|W|}\max_{i\in W}z_i$, and sorting
$z_{(1)}\ge\dots\ge z_{(n)}\ge z_{(n+1)}=0$ gives
\[
\sum_WR_{|W|}\max_{i\in W}z_i=\sum_{k=1}^{n}\bigl(z_{(k)}-z_{(k+1)}\bigr)C_k .
\]
If $g$ satisfies every subset condition, $z\cdot g$ is at most this value
for every $z\ge0$. The nonnegative vectors dominated by resource allocations
form a compact downward-closed polytope; separation therefore gives flows
$f_{W,i}\ge0$ with $\sum_{i\in W}f_{W,i}\le R_{|W|}$ and
$\sum_{W\ni i}f_{W,i}\ge g_i$. We distribute unused resources within each $W$
and pay its members $f_{W,i}/(2c_{|W|})$, with zero to others, on both paired
profiles. For $|W|=m+1$, we instead pay each majority member $V/(m+1)$.
These payments are label-symmetric and exhaust every profile's budget;
minimal majorities contribute zero premium, so
\[
G_i(v)=\sum_{W\ni i}f_{W,i}\ge g_i .\qedhere
\]
\end{proof}

\paragraph{One validator with at least half the stake.} Let $n=3$, let
validator $1$ hold stake $\frac12\le\varphi_1<1$ and the other two equal
positive stakes, with common $\eta$ and $\omega$, and write $t=pq$. Averaging over the two small
validators keeps quality and every constraint, so an allocation is described
by $0\le a_0\le a_1\le a_2\le\frac12$, miner $1$'s share when validator $1$
reports $0$ and none, one or both small validators report $1$; the other
profiles follow by symmetry, and at $\varphi_1=\frac12$ ties force
$a_2=\frac12$. Integrating over profiles,
\begin{gather*}
Q=p+d\,[\,ta_2-2ta_1-(1-t)a_0\,],\\
H_1=1-(1-2t)(a_0+a_2)-4ta_1,\\
H_S=(1-4t)a_1-(1-2t)a_0+2ta_2,
\end{gather*}
where $H_1$ and $H_S$ are the influences of the large and of a small validator.
With three validators only unanimous profiles create premium, so every
nonempty capacity equals $C_k=Vd^2/2$ and the subset conditions of
Theorem~\ref{thm:capacity} reduce to
$3\eta+\frac12\omega M(H_1+2H_S)\le Vd^2/2$. Thus no rule is feasible
if $\eta>Vd^2/6$; we now assume $0\le\eta\le Vd^2/6$.
For $\varphi_1>1/2$, the allocation polytope has vertices $(\frac12,\frac12,\frac12)$, $(0,0,0)$, $(0,0,\frac12)$ and
$(0,\frac12,\frac12)$, whose pairs (quality gain over one half divided by $d$,
total influence $H_1+2H_S$) are $(0,0)$, $(\frac12,1)$,
$(\frac{1+t}2,\frac12+3t)$ and $(\frac{1-t}2,\frac32-3t)$. The third vertex has
the highest quality gain and the highest ratio of gain to influence; its cross
products with the second and fourth differ by nonnegative multiples of $1-4t$.
Mixing it with the even split is therefore optimal. Both vertices lie on
the face $a_2=1/2$, so the same mixture is feasible and optimal at
$\varphi_1=1/2$. If $\omega=0$, its undamped quality is optimal for every
feasible cost. For $\omega>0$, the constraint gives
\[
Q^*=\frac12+\frac{d(1+t)}2\min\Bigl\{1,\frac{Vd^2-6\eta}{\omega M(\frac12+3t)}\Bigr\}.
\]

\begin{proof}[Proof of Theorem~\ref{thm:memory}]
We first bound the pre-acquisition belief $\pi$. Each round reveals at most
$n$ fresh signals, giving likelihood ratio in $[\Lambda^{-1},\Lambda]$.
The Markov transition confines $\pi$ to $[1-\Pi,\Pi]$; unanimous honest
histories approach both endpoints. A blind report $1$ has affine validation
pay $W_1(\pi)$, with $W_1(1/2)=V/n-G(v)$ and $W_1(p)=V/n$ by symmetry.
Its miner-allocation gain is $(p-d\pi)H(a)$, giving
$2(p-\pi)G(v)/d\ge\eta+\omega M(p-d\pi)H(a)$.
The label-flipped constraint covers report $0$. Blind-report losses equal
probability-weighted conditional misreport losses, so neither acquired
misreports nor mixtures gain.

For fixed $\beta_i$, honest flow payoff is
$V/n-\eta+\beta_iM[1-Q+(2Q-1)\pi]$, $Q=Q(a)$.
Iterating affine Markov prediction and discounting therefore gives an
affine bounded honest continuation value $U(\pi)$.
For every current policy, $\E[\pi'\mid\text{history}]=(1-\rho)/2+\rho\pi$,
so affine $U$ makes expected continuation independent of acquisition and
reporting. The current constraints give
$\E[u_i+\delta U(\pi')\mid\text{history}]\le U(\pi)$.
Bounded $U$ lets this telescope for arbitrary randomized history-dependent
deviations. A strict endpoint violation yields a profitable one-round
deviation at a finite unanimous history, followed by honest play.

Feasibility implies $G(v)\ge\eta>0$. The ratio
$[2(p-\pi)G(v)/d-\eta]/(p-d\pi)$ has derivative
$-(4pqG(v)/d+d\eta)/(p-d\pi)^2<0$, so $\Pi$ binds.
If $\Pi\ge p$, positive cost is infeasible. Otherwise
Lemmas~\ref{lem:rent} and~\ref{lem:influence} give the frontier, attained by
damped majority and majority sharing; at $\omega=0$, full majority is feasible
exactly when $\eta\le G_\rho$. Setting $\Pi=p$ in its fixed-point equation gives
$\rho_c=d(p^{n+1}+q^{n+1})/(p^{n+1}-q^{n+1})$.\qedhere
\end{proof}

\section{Replication details}
\label{app:replica}

\paragraph{Implementation.} The port fixes Subtensor commit c004cebf and
takes effective stake and resolved chain state; it does not rebuild the
runtime. Sparse zeros are retained for adaptive bonds. Additional tests
comprise 10 classic epochs, 113 further epochs and 23 helper inputs.
Some official Yuma3 bond assertions iterate over empty ranges; the
differential comparison checks every bond and seven intermediate traces.

\begin{table}[t]
\caption{Closed-form ownership tolerance versus caps from sampled
16-bit deviations (Theorem~\ref{thm:hedging} and
Proposition~\ref{prop:soft}). Running example, bond memory off.}
\label{tab:replica}
\small
\begin{tabular}{@{}rrrr@{}}
\toprule
$\ell$ & closed form $\tau(\ell)$ & sampled cap & quality \\
\midrule
$0$ & $0.009259$ & $0.009275$ & $0.843750$ \\
$0.1$ & $0.053782$ & $0.053778$ & $0.764767$ \\
$0.2$ & $0.088921$ & $0.088909$ & $0.692188$ \\
$0.4$ & $0.139172$ & $0.139149$ & $0.560714$ \\
$0.49$ & $0.155275$ & $0.155262$ & $0.505951$ \\
\bottomrule
\end{tabular}
\end{table}

\paragraph{Threshold calculations.} Table~\ref{tab:replica} samples
fixed-sum $65535$ reports: about 1040 deviations per $\ell$, including
nearby boundaries, yield 15,657 epochs and 20 additional Rust comparisons.
The encoded $\ell=0.1,0.49$ are $0.1000076,0.4899977$.
These sampled caps do not certify all integer reports.
For Table~\ref{tab:thresholds}, the data retain exact rational values for
all odd $n=3,\ldots,51$ at the three displayed accuracies.

\paragraph{Single-validator experiment.} The grid uses $n=11$ with
$J=2,5,10$, and $n=31$ with $J=5$, each at $s_0=1/n,0.1,0.25,0.4$.
Both $h=0.1,0.25$ are fixed in advance. All-or-nothing controls use
$n=11$, $s_0=0.4$, $J=2,5$. Miners have separate identities and zero
stake. Stakes and epoch emission use scale $10^9$; each weight row is
scaled to maximum $65535$ and rounded.

The seed is $20260930+100000n+1000J$. Reports and stake variants share
signals. The 2048 paired draws include the 512-draw pilot as a prefix;
intervals use paired standard errors and $t_{2047}$ critical values,
without multiple-comparison adjustment. Eighteen distinct new inputs
match all Rust outputs and traces. All 36 comparisons are in the data.

\paragraph{Fallback.} If every miner column is clipped to zero, the source
pays the entire emission to validators. This occurs in 31.15\% of
unhedged five-miner control draws and none after hedging; at $h=0.1$,
income falls by $0.05273V$, using the fixed nominal budget $V$.
The two-miner controls' income intervals contain zero.
No fallback occurs in the softmax sample. Quality is set to zero when
miner emission vanishes; the data also report best-miner income
relative to total emission.
Replication package: \href{https://github.com/submissionrepo/bittensor-yuma-replication}{GitHub}.
The package includes exact/symbolic checks and independent LLM audits.
\end{document}
